% !TeX program = lualatex
% !TeX encoding = utf8
% !TeX spellcheck = uk_UA
% !TeX root =../PyscfBook.tex


%% ========================================================
\section{Метод зв’язаних кластерів}
%% ========================================================

Метод зв'язаних кластерів (\texttt{Coupled Cluster, CC}) є природним кроком після методу
конфігураційної взаємодії (\texttt{CI}), але з принципово новим підходом
до побудови хвильової функції.

У методі \texttt{CI} хвильова функція системи записується як лінійна комбінація
детермінантів, збуджених відносно основного (\texttt{HF}) детермінанта \eqref{eq:CI}.
Така форма зручна для чисельної реалізації, проте має два
фундаментальні недоліки:
\begin{enumerate}
    \item Щоб отримати точний результат, потрібно включити всі можливі
    збудження (\texttt{Full CI}), що обчислювально неможливо для більшості систем.
    \item Хвильова функція \texttt{CI} не є \textbf{size-extensive}, тобто енергія двох незалежних систем
    не дорівнює сумі їхніх енергій. Це призводить до помилок при описі
    великих молекул.
\end{enumerate}

Метод зв'язаних кластерів (\texttt{CC}) розв'язує ці проблеми, вводячи іншу,
експоненціальну форму хвильової функції:
\begin{equation}\label{eq:CC}
    |\Psi_{CC}\rangle = e^{\hat{T}} |\Phi_0\rangle
\end{equation}
де $|\Phi_0\rangle$ --- HF-детермінант, $\hat{T}$ --- \textbf{кластерний оператор} $\hat{T}$:
\begin{equation}
    \hat{T} = \hat{T}_1 + \hat{T}_2 + \hat{T}_3 + \cdots,
\end{equation}
а кожен член \(\hat{T}_n\) відповідає всім можливим \(n\)-кратним збудженням
відносно детермінанта Хартрі–Фока $|\Phi_0\rangle$:
\begin{align}
    \hat{T}_1 &= \sum_{i,a} t_i^a\, a_a^\dagger a_i, \\
    \hat{T}_2 &= \frac{1}{4} \sum_{i,j,a,b} t_{ij}^{ab}\, a_a^\dagger a_b^\dagger a_j a_i, \\
    &\vdots
\end{align}

На перший погляд, експоненціальна форма виглядає складнішою, ніж
звичайна лінійна комбінація. Проте її фундаментальна перевага полягає
у властивості експоненти:
\begin{equation}
    e^{\hat{T}} = 1 + \hat{T} + \frac{1}{2!}\hat{T}^2 +
    \frac{1}{3!}\hat{T}^3 + \cdots,
\end{equation}
що означає, що навіть якщо явно враховуються лише \(\hat{T}_1\) і \(\hat{T}_2\)
(метод \texttt{CCSD}), хвильова функція автоматично містить внески від
трикратних, чотирикратних та вищих збуджень через добутки
\(\hat{T}_1^2\), \(\hat{T}_1\hat{T}_2\) тощо.

%Таким чином, \texttt{CC} дозволяє врахувати ефекти високих порядків
%кореляції, використовуючи відносно малу кількість параметрів
%(\(t_i^a,\, t_{ij}^{ab}, \ldots\)).
%Крім того, енергія, отримана в методі \texttt{CC}, є \textbf{size-extensive}, тобто правильно
%масштабується з розміром системи --- необхідна властивість для опису
%великих молекул і конденсованих систем.

%% --------------------------------------------------------
\subsection{Виведення рівнянь методу CC}
%% --------------------------------------------------------

Почнемо з основного рівняння Шредінґера для хвильової функції
методу зв'язаних кластерів:
\begin{equation}
    \hat{H} \Psi_{\text{CC}} = E_{\text{CC}} \Psi_{\text{CC}},
\end{equation}
де \(\Psi_{\text{CC}} = e^{\hat{T}} \Phi_0\), а \(\Phi_0\) --- детермінант Хартрі–Фока.

Щоб отримати рівняння для невідомих амплітуд \(t_i^a, t_{ij}^{ab}, \ldots\),
підставимо цю форму у рівняння Шредінґера та ліворуч домножимо на
\(e^{-\hat{T}}\):
\begin{equation}
    e^{-\hat{T}} \hat{H} e^{\hat{T}} \Phi_0 = E_{\text{CC}} \Phi_0.
\end{equation}

Позначимо ефективний гамільтоніан як
\begin{equation}
    \bar{H} = e^{-\hat{T}} \hat{H} e^{\hat{T}}.
\end{equation}
Тоді рівняння Шредінґера набуває вигляду
\[
\bar{H} \Phi_0 = E_{\text{CC}} \Phi_0.
\]

Щоб знайти енергію, проєктуємо це рівняння на основний детермінант:
\begin{equation}
    E_{\text{CC}} = \langle \Phi_0 | \bar{H} | \Phi_0 \rangle.
\end{equation}
Це рівняння визначає енергію системи.

Далі, щоб знайти коефіцієнти \(t_i^a, t_{ij}^{ab}\), проєктуємо те саме рівняння
на збуджені детермінанти \(\Phi_i^a, \Phi_{ij}^{ab}, \ldots\):
\begin{equation}
    \langle \Phi_i^a | \bar{H} | \Phi_0 \rangle = 0, \qquad
    \langle \Phi_{ij}^{ab} | \bar{H} | \Phi_0 \rangle = 0, \quad \text{і т. д.}
\end{equation}
Тобто ми вимагаємо, щоб ефективний гамільтоніан \(\bar{H}\)
не створював збуджень із основного стану --- лише тоді
експоненціальна хвильова функція \(e^{\hat{T}} \Phi_0\) буде власним станом
оператора \(\hat{H}\).

\bigskip

\textbf{Фізичний зміст.}
Експоненціальний оператор \(e^{\hat{T}}\) створює всі можливі збудження
основного детермінанта, але в такий спосіб, що кожне з них враховується
один раз і з правильним знаком. Тому метод автоматично виключає
``повторні'' або незалежні частини збуджень, залишаючи лише ті,
що дійсно взаємопов'язані між електронами. Саме тому метод називається
\emph{метод зв'язаних кластерів}.



%% --------------------------------------------------------
\subsection{Доведення \emph{size-extensivity}}
%% --------------------------------------------------------

Нехай система складається з двох несв'язаних підсистем \(A\) та \(B\),
які не взаємодіють між собою (відстань між ними → \(\infty\) або відсутні міжсистемні члени в гамільтоніані).
Позначимо гамільтоніан повної системи:
\[
\hat{H} = \hat{H}_A + \hat{H}_B,
\]
і виберемо базовий детермінант як простий тензорний добуток:
\[
\Phi_0 = \Phi_0^A \otimes \Phi_0^B.
\]

У методі \texttt{CC} кластерний оператор для повної системи при відсутності взаємодії розкладається в суму операторів, що діють у відповідних підпросторах:
\[
\hat{T} = \hat{T}_A + \hat{T}_B,
\]
оскільки збудження, що зачіпають одночасно і \(A\), і \(B\), відсутні.

Тоді експоненційний анзац факторизується:
\[
e^{\hat{T}} = e^{\hat{T}_A + \hat{T}_B} = e^{\hat{T}_A}\, e^{\hat{T}_B},
\]
оскільки \([\hat{T}_A,\hat{T}_B]=0\) (оператори діють у різних підпросторах).

Розглянемо ефективний гамільтоніан
\[
\bar{H} = e^{-\hat{T}} \hat{H} e^{\hat{T}}.
\]
Підставляючи розклади для \(\hat{H}\) і \(\hat{T}\) отримуємо:
\[
\bar{H} = e^{-(\hat{T}_A+\hat{T}_B)}(\hat{H}_A+\hat{H}_B)e^{\hat{T}_A+\hat{T}_B}
= e^{-\hat{T}_A}\hat{H}_A e^{\hat{T}_A} + e^{-\hat{T}_B}\hat{H}_B e^{\hat{T}_B},
\]
оскільки, наприклад, \(e^{-\hat{T}_A}\hat{H}_B e^{\hat{T}_A}=\hat{H}_B\) (оператори \(\hat{H}_B\) і \(\hat{T}_A\) комутують)
і аналогічно для інших перехресних членів.

Енергія \texttt{CC} визначається як проєкція на основний детермінант:
\[
E_{\text{CC}} = \langle \Phi_0 | \bar{H} | \Phi_0 \rangle.
\]
Підставляючи попередній розклад \(\bar{H}\) та використовуючи тензорну структуру \(\Phi_0\), маємо:
\[
\begin{aligned}
E_{\text{CC}}
&= \langle \Phi_0^A\otimes\Phi_0^B |
\bigl(e^{-\hat{T}_A}\hat{H}_A e^{\hat{T}_A} + e^{-\hat{T}_B}\hat{H}_B e^{\hat{T}_B}\bigr)
| \Phi_0^A\otimes\Phi_0^B \rangle \\
&= \langle \Phi_0^A | e^{-\hat{T}_A}\hat{H}_A e^{\hat{T}_A} | \Phi_0^A \rangle
\;+\;
\langle \Phi_0^B | e^{-\hat{T}_B}\hat{H}_B e^{\hat{T}_B} | \Phi_0^B \rangle \\
&= E_{\text{CC}}(A) + E_{\text{CC}}(B).
\end{aligned}
\]
Отже, енергія методу \texttt{CC} для несв'язаних підсистем є адитивною,
тобто метод є \emph{size\-extensive}.

\bigskip

\begin{commentbox}
У лінійному анзацi \texttt{CI} хвильова функція є сумою детермінантів,
і при розбитті на підсистеми з'являються незв'язані (розкладні) члени,
які приводять до крос-термінів у енергії. Через це лінійний \texttt{CI}
(за винятком \texttt{Full CI}) не гарантує адитивності енергії і тому не є
size-extensive. Експоненціальний анзац \(\Psi_{\text{CC}}=e^{\hat{T}}\Phi_0\)
автоматично виключає такі незв'язані повтори і забезпечує правильну
адитивну поведінку енергії.
\end{commentbox}

%% --------------------------------------------------------
\subsection{Наближення в методі CC: \texttt{CCSD} і \texttt{CCSD(T)}}
%% --------------------------------------------------------

У практичних розрахунках повна експоненційна форма хвильової функції методу
зв'язаних кластерів:
\[
|\Psi_{\text{CC}}\rangle = e^{\hat{T}} |\Phi_0\rangle,
\]
де
\[
\hat{T} = \hat{T}_1 + \hat{T}_2 + \hat{T}_3 + \dots + \hat{T}_N,
\]
залишається, звісно, нескінченною, адже кількість можливих багатоелектронних
збуджень зростає комбінаційно з числом електронів.
Тому необхідно робити \textbf{наближення}, обмежуючи порядок збуджень.

\paragraph{Наближення \texttt{CCSD}}
Найпоширеніше робоче наближення — це
\[
\hat{T} \approx \hat{T}_1 + \hat{T}_2,
\]
тобто враховуються лише \textbf{однократні} та \textbf{дворазові збудження}.
Таке наближення позначається як \texttt{CCSD}
(від англ. \emph{Coupled Cluster with Single and Double excitations}).

Оператор хвильової функції має вигляд:
\[
|\Psi_{\text{CCSD}}\rangle = e^{\hat{T}_1 + \hat{T}_2} |\Phi_0\rangle.
\]

Хоча формально ми не включаємо потрійні й вищі збудження,
експоненційна форма гарантує, що частина їхнього ефекту враховується автоматично
через добутки нижчих операторів, наприклад:
\[
\frac{1}{2!}\hat{T}_1^2 \sim \text{подвійні збудження}, \quad
\hat{T}_1 \hat{T}_2 \sim \text{потрійні збудження}.
\]
Тобто навіть \texttt{CCSD} описує взаємодію між різними рівнями кореляції значно точніше,
ніж метод \texttt{CISD} (конфігураційна взаємодія з одиничними та подвійними збудженнями).

\paragraph{Корекція \texttt{(T)}: «золотий стандарт»}
Для більшої точності часто додають внесок потрійних збуджень \(\hat{T}_3\)
у вигляді \emph{неітераційної поправки}, яку обчислюють на основі вже отриманих
амплітуд \(\hat{T}_1, \hat{T}_2\).
Таке наближення позначається як \texttt{CCSD(T)}.

Символ \texttt{(T)} означає, що потрійні збудження враховані
\emph{perturbatively}, тобто за теорією збурень.
Така схема зберігає високу точність, наближену до повного
\texttt{CCSDT}, але з обчислювальною складністю всього
\(\mathcal{O}(N^7)\) замість \(\mathcal{O}(N^8)\)–\(\mathcal{O}(N^{10})\)
для повного варіанта.

\paragraph{Практичне значення}
Метод \texttt{CCSD(T)} виявився настільки надійним у відтворенні
експериментальних енергій, геометрій і частот коливань, що його часто називають
\textbf{«золотим стандартом» квантової хімії}.
Для невеликих і середніх молекул похибка енергії зв’язку зазвичай не перевищує
\(1{-}2~\text{kcal/mol}\), тобто досягає хімічної точності.

\begin{center}
\begin{tblr}{
colspec={Q[l,m] Q[l,m] Q[l,m]},
row{1} = {c,m, font=\bfseries}
}
\toprule
Метод & Враховані збудження & Тип наближення \\ \midrule
\texttt{CCSD} & \(\hat{T}_1, \hat{T}_2\) & Ітераційне \\
\texttt{CCSD(T)} & \(\hat{T}_1, \hat{T}_2\) + корекція \(\hat{T}_3\) & Ітераційно + збурювально \\
\texttt{CCSDT} & \(\hat{T}_1, \hat{T}_2, \hat{T}_3\) & Повне, ітераційне \\ \bottomrule
\end{tblr}
\end{center}

Таким чином, \texttt{CCSD(T)} поєднує фізичну повноту опису з практичною
ефективністю й залишається основним методом високоточної
квантової хімії для систем до кількох десятків атомів.


%%% --------------------------------------------------------
%\subsubsection{Доведення size-extensive для Coupled Cluster}
%%% --------------------------------------------------------
%
%Розглянемо дві невзаємодіючі підсистеми \(A\) і \(B\), для яких
%гамільтоніан має вигляд:
%\[
%\hat{H} = \hat{H}_A + \hat{H}_B.
%\]
%Основний (референсний) детермінант системи є добутком детермінантів підсистем:
%\[
%|\Phi_0\rangle = |\Phi_A\rangle \otimes |\Phi_B\rangle.
%\]
%
%У методі Coupled Cluster хвильова функція має експоненціальну форму:
%\[
%|\Psi_{CC}\rangle = e^{\hat{T}}|\Phi_0\rangle,
%\]
%де кластерний оператор розкладається як
%\[
%\hat{T} = \hat{T}_A + \hat{T}_B.
%\]
%Оскільки \(\hat{T}_A\) і \(\hat{T}_B\) діють у різних підпросторах (на різні набори орбіталей),
%вони комутують:
%\[
%[\hat{T}_A, \hat{T}_B] = 0.
%\]
%Отже, експонента факторизується:
%\[
%e^{\hat{T}} = e^{\hat{T}_A+\hat{T}_B} = e^{\hat{T}_A} e^{\hat{T}_B}.
%\]
%
%Енергія CC визначається як середнє значення подібно перетвореного гамільтоніана:
%\[
%E_{\mathrm{CC}} = \langle \Phi_0 | \bar{H} | \Phi_0 \rangle,
%\qquad
%\bar{H} = e^{-\hat{T}} \hat{H} e^{\hat{T}}.
%\]
%Підставляючи \(\hat{H}=\hat{H}_A+\hat{H}_B\) і \(\hat{T}=\hat{T}_A+\hat{T}_B\), маємо:
%\[
%\bar{H}
%= e^{-(\hat{T}_A+\hat{T}_B)} (\hat{H}_A+\hat{H}_B) e^{\hat{T}_A+\hat{T}_B}
%= e^{-\hat{T}_A}\hat{H}_A e^{\hat{T}_A}
%+ e^{-\hat{T}_B}\hat{H}_B e^{\hat{T}_B},
%\]
%оскільки перехресні члени, що містять одночасно \(\hat{H}_A\) та \(\hat{T}_B\),
%зникають (оператори діють у різних підпросторах і комутують).
%
%Отже, повна енергія системи розкладається як:
%\[
%E_{\mathrm{CC}} =
%\langle \Phi_A \Phi_B |
%\big(e^{-\hat{T}_A}\hat{H}_A e^{\hat{T}_A} + e^{-\hat{T}_B}\hat{H}_B e^{\hat{T}_B}\big)
%| \Phi_A \Phi_B \rangle
%= E_{\mathrm{CC}}(A) + E_{\mathrm{CC}}(B).
%\]



%\begin{commentbox}
%Розмірна узгодженість (size-consistency) є необхідною властивістю для коректного опису дисоціаційної
%границі молекули. Однак навіть розмірно узгоджені методи (наприклад, CCSD) не гарантують
%фізично правильної поведінки потенціальної кривої при розриві зв’язку.
%Причина полягає у руйнуванні однодетермінантного опису --- хвильова функція стає
%багатоконфігураційною, а референс Хартрі–Фока перестає бути адекватним.
%У таких випадках необхідно застосовувати мультидетермінантні методи (CASSCF, MRCI),
%або використовувати UCCSD для часткової корекції ефектів спінового розшарування.
%\end{commentbox}


%\paragraph{CCSD (Coupled Cluster Singles and Doubles).}
%Враховує лише одно- та двозбуджені конфігурації:
%\begin{align}
%    \hat{T}_1 & = \sum_{i}^{\text{occ}} \sum_{a}^{\text{virt}} t_i^a \, \hat{a}^\dagger_a \hat{a}_i \\
%    \hat{T}_2 & = \frac{1}{4} \sum_{ij}^{\text{occ}} \sum_{ab}^{\text{virt}} t_{ij}^{ab} \, \hat{a}^\dagger_a \hat{a}^\dagger_b \hat{a}_j \hat{a}_i
%\end{align}
%
%\paragraph{CCSD(T) --- «золотий стандарт».}
%Для ще точнішого врахування кореляції додають пертурбативну поправку від потрійних збуджень:
%\begin{equation}
%    E_{CCSD(T)} = E_{CCSD} + E_{(T)}
%\end{equation}
%де $E_{(T)}$ визначається за формулами четвертого порядку теорії збурень (MBPT4).
%Метод CCSD(T) поєднує високу точність з помірною обчислювальною складністю $O(N^7)$ і тому
%вважається \textbf{еталоном точності} у квантовій хімії.
%
%
%Експоненційна форма хвильової функції забезпечує:
%\begin{itemize}
%    \item \textbf{Size-extensivity:} на відміну від CISD, енергія двох незалежних систем у CCSD дорівнює сумі енергій кожної з них;
%    \item \textbf{Автоматичне включення ефективних багатоелектронних збуджень:} хоча $\hat{T}$ містить лише $T_1$ і $T_2$, експонента $e^{\hat{T}}$ генерує вищі збудження ($T_3, T_4, \dots$);
%    \item \textbf{Високу стабільність} при обчисленнях навіть для систем із помірним багатоконфігураційним характером.
%\end{itemize}


%%% --------------------------------------------------------
%\subsection{CCSD розрахунок атома Гелію}
%%% --------------------------------------------------------
%
%Система гелію (\ce{He}) --- це найпростіший приклад, де метод
%\texttt{CCSD} дає \textbf{точний результат у межах будь-якого базису}.
%Причина полягає в тому, що для двоелектронних систем усі кореляційні ефекти
%повністю описуються лише оператором парних збуджень \(T_2\),
%тоді як оператор однократних збуджень \(T_1\) зникає через симетрію.
%
%\paragraph{Фізична суть.}
%
%Для атома гелію існує лише одна заповнена орбіталь \(1s^2\).
%Будь-яке корельоване збудження полягає в тому, що обидва електрони
%одночасно переходять у віртуальні орбіталі:
%\[
%T_2 = \sum_{abij} t_{ij}^{ab}\, a_a^\dagger a_b^\dagger a_j a_i.
%\]
%Оператор \(T_1\) відсутній, оскільки немає однократних збуджень,
%які не порушували б симетрію стану.
%
%Таким чином, хвильова функція CCSD спрощується до:
%\[
%|\Psi_{\text{CCSD}}\rangle = e^{T_2} |\Phi_0\rangle,
%\]
%що збігається з точним розв’язком \texttt{Full CI} у даному базисі:
%\[
%E_{\text{CCSD}} = E_{\text{FCI}}.
%\]
%
%\paragraph{Методичне значення прикладу.}
%
%Розрахунок гелію є корисним із кількох причин:
%\begin{itemize}
%  \item Він служить \textbf{еталонним тестом точності} реалізації CCSD у будь-якій програмі:
%        результат має збігатися з Full CI до \(10^{-10}\,E_h\).
%  \item Дає змогу продемонструвати, що для систем із двома електронами
%        вся електронна кореляція є \emph{парною}, тобто описується лише $T_2$-оператором.
%  \item Дозволяє спостерігати \textbf{зникнення однократних збуджень} $T_1$ як наслідок симетрії,
%        що особливо важливо для розуміння поведінки CCSD у багатоелектронних системах.
%  \item Є \textbf{контрольним прикладом збіжності}:
%        ітераційна процедура CCSD для гелію збігається за 2–3 кроки.
%\end{itemize}
%
%\paragraph{Фізичний результат.}
%
%Для гелію різниця між енергіями:
%\[
%E_{\mathrm{HF}}, \quad E_{\mathrm{CCSD}}, \quad E_{\mathrm{FCI}}
%\]
%становить лише кілька мілігаусів (mHa), тобто CCSD повністю відтворює кореляційну енергію.
%Це підтверджує, що метод кластерів є \emph{систематичним і точним наближенням},
%яке починає відхилятись від FCI лише тоді, коли з’являються багаточастинкові
%(потрійні, четверні) кореляційні ефекти.
%
%\paragraph{Висновок.}
%
%Цей розрахунок не є практично «новим результатом», але є
%\textbf{методично показовим}:
%він демонструє, що CCSD --- це узагальнення ідей Хартрі–Фока,
%яке в межах двоелектронної системи дає точну відповідь,
%а для складніших систем --- стає точним першим наближенням до повного опису кореляції.
%
%Прклад розрахунку подано у файлі \inputcode{He_CCSD.py}
%
%%% --------------------------------------------------------
%\subsection{CCSD для відкритих оболонок (UCCSD)}
%%% --------------------------------------------------------
%
%\inputcode{code_9.py}
%
%Для систем із неспареними електронами використовується варіант \textbf{UCCSD (Unrestricted CCSD)}.
%Він будується на UHF-детермінанті, дозволяючи різні орбіталі для $\alpha$- та $\beta$-спінів.
%UCCSD зберігає всі основні переваги CCSD і забезпечує адекватний опис навіть для радикалів та іонів.
%
%%% --------------------------------------------------------
%\subsection{Перевірка size-extensivity: дисоціація \ce{LiH}}
%%% --------------------------------------------------------
%
%\inputcode[mode=inline]{LiH_size_consistency.py}
%
%Порівняння з CISD демонструє ключову відмінність CCSD:
%для повністю розірваної молекули \ce{LiH} (великі між’ядерні відстані)
%енергія \textbf{CISD} не дорівнює сумі енергій окремих атомів (\textit{порушення size-consistency}),
%тоді як \textbf{CCSD} дає правильний результат:
%\[
%E_{\text{CCSD}}(\ce{LiH}_\text{diss}) \approx E_{\text{CCSD}}(\ce{Li}) + E_{\text{CCSD}}(\ce{H})
%\]
%Це одна з найважливіших переваг CC над CI при моделюванні хімічних реакцій і дисоціацій.
%


